\documentclass[fleqn]{beamer}

\mode<presentation> {
\usetheme{Madrid}
\usecolortheme{dolphin}
}

\usepackage[utf8]{inputenc}
\usepackage[brazilian]{babel}
\usepackage{amsmath}
\usepackage{amssymb}
\usepackage{booktabs}
\usepackage{listings}
\usepackage{xcolor}

\lstdefinelanguage{AMPL}{
	morekeywords={param, var, subject, to, maximize, minimize, model, data,
		option, solve, display, solver, sum, in, binary, symbolic, Infinity, printf, for, if, then},
	sensitive=true,
	morecomment=[l]{\#},
	morestring=[b]",
}
\lstset{
	basicstyle=\ttfamily\scriptsize,
	keywordstyle=\color{blue!70!black}\bfseries,
	commentstyle=\color{gray},
	stringstyle=\color{orange!70!black},
	breaklines=true,
	frame=single,
	backgroundcolor=\color{black!3},
}

\title[Aula 1 -- Introdução à Otimização]{Introdução à Otimização}
\subtitle{Modelagem, AMPL e o Problema da Mochila}
\author[Santi]{Prof. \'Everton Santi}
\institute[UFRN/ECT]{ECT2524 -- Introdução à Otimização \\ Escola de Ciências e Tecnologia -- UFRN}
\date{Aula 1}

\begin{document}

\begin{frame}
	\titlepage
\end{frame}

\begin{frame}
	\frametitle{Sumário}
	\tableofcontents
\end{frame}

%------------------------------------------------------------
\section{O que é otimização?}
%------------------------------------------------------------

\begin{frame}
	\frametitle{O que significa otimização?}
	Um problema de otimização pode ser expresso como:
	$$
		\text{min}~f(x)
	$$
	sujeito a
	$$
		x \in S
	$$
	em que:
	\begin{itemize}
		\item $x$ é a \textbf{solução} do problema (variáveis de decisão);
		\item $f(x)$ é a \textbf{função objetivo};
		\item $S$ é o \textbf{conjunto de soluções viáveis}.
	\end{itemize}
\end{frame}

\begin{frame}
	\frametitle{Função objetivo -- exemplos}
	A função objetivo pode representar, por exemplo:
	\begin{itemize}
		\item Distância para se deslocar de uma cidade até outra;
		\item Desperdício de material na fabricação de um item;
		\item Custo de uma refeição que atenda exigências nutricionais;
		\item Tempo total de execução de tarefas;
		\item Custo de construção de uma rede de distribuição de água;
		\item Entre muitos outros...
	\end{itemize}
\end{frame}

\begin{frame}
	\frametitle{Restrito vs. irrestrito}
	O conjunto de soluções viáveis $S$ é delimitado por um \textbf{conjunto de
	restrições}.
	\begin{itemize}
		\item Se $S = \emptyset$ (sem restrições declaradas): problema
			\textbf{irrestrito};
		\item Se $S \neq \emptyset$: problema \textbf{restrito}.
	\end{itemize}
	\vspace{0.5cm}
	\textit{Nota}: um problema restrito ainda pode ter infinitas soluções
	viáveis -- restrito não significa finito.
\end{frame}

%------------------------------------------------------------
\section{Modelagem}
%------------------------------------------------------------

\begin{frame}
	\frametitle{Modelagem}
	Qualquer situação do mundo real que possa ser representada por
	expressões matemáticas pode ser tratada com otimização, desde que o
	modelo permita uma aproximação satisfatória da realidade.
	\vspace{0.3cm}

	Um \textbf{modelo} é uma representação do mundo real. Em otimização,
	falamos de um \textbf{modelo simbólico}: parâmetros, variáveis,
	índices, funções, igualdades e desigualdades.
	\vspace{0.3cm}

	Modelar é um misto de \textbf{arte e ciência}: exige criatividade para
	observar, entender e traduzir a realidade, além de conhecimento técnico
	para tratá-la adequadamente.
\end{frame}

%------------------------------------------------------------
\section{Ferramentas computacionais}
%------------------------------------------------------------

\begin{frame}
	\frametitle{Linguagem de modelagem e solver}
	\begin{itemize}
		\item \textbf{Linguagem de modelagem}: escrever a formulação de
			forma próxima da notação matemática (AMPL, GAMS, ZIMPL, OPL, ...);
		\item \textbf{Resolvedor (\textit{solver})}: programa
			especializado que, dados os parâmetros e a formulação, busca
			os valores das variáveis de decisão (CPLEX, Gurobi, HiGHS,
			CBC, ...).
	\end{itemize}
	\vspace{0.4cm}
	Neste curso usaremos \textbf{AMPL}: linguagem simples de aprender,
	com versão gratuita, e que permite usar diversos \textit{solvers}.
\end{frame}

\begin{frame}
	\frametitle{Estrutura básica de um projeto AMPL}
	\begin{itemize}
		\item \texttt{.mod} -- formulação do modelo (parâmetros, variáveis,
			objetivo, restrições);
		\item \texttt{.dat} -- dados de entrada (valores dos parâmetros);
		\item \texttt{.run} -- comandos para carregar modelo + dados e
			resolver.
	\end{itemize}
	\vspace{0.4cm}
	\textit{Nota}: instruções AMPL terminam em \texttt{;} -- diferente de
	Python, que não usa isso. Atenção ao alternar entre as duas linguagens
	no Colab.
\end{frame}

%------------------------------------------------------------
\section{Problema da Mochila}
%------------------------------------------------------------

\begin{frame}
	\frametitle{Um ``Hello World'' diferente}
	Você vai acampar no fim de semana e encontra uma mochila com
	capacidade de \textbf{10Kg} e 6 itens candidatos, cada um com peso e
	uma pontuação de importância (1 a 10):
	\begin{itemize}
		\item Item 1: Lanterna (1Kg, 8 pontos)
		\item Item 2: Barraca (4Kg, 10 pontos)
		\item Item 3: Reservatório de água (3Kg, 10 pontos)
		\item Item 4: Reservatório de água (2Kg, 8 pontos)
		\item Item 5: Colchão inflável (0.5Kg, 7 pontos)
		\item Item 6: Caixa de som (2Kg, 4 pontos)
	\end{itemize}
	Decida quais itens levar, \textbf{maximizando} a soma das pontuações
	sem \textbf{ultrapassar} a capacidade da mochila.
\end{frame}

\begin{frame}
	\frametitle{Formulação expandida}
	$$
		\text{maximize}~Z = 8x_1 + 10x_2 + 10x_3 + 8x_4 + 7x_5 + 4x_6
	$$
	sujeito a:
	$$
		x_1 + 4x_2 + 3x_3 + 2x_4 + 0.5x_5 + 2x_6 \leq 10
	$$
	$$
		x_1, x_2, x_3, x_4, x_5, x_6 \in \{0,1\}
	$$
	\vspace{0.3cm}
	\textbf{Limitação}: só serve para esta instância específica -- 6
	itens fixos, escrito ``na mão''.
\end{frame}

\begin{frame}
	\frametitle{Formulação compacta}
	Parâmetros:
	\begin{itemize}
		\item $n$ -- total de itens disponíveis;
		\item $C$ -- capacidade de carga da mochila;
		\item $p_i$ -- peso do item $i$, $\forall i = 1,\dots,n$;
		\item $b_i$ -- benefício do item $i$, $\forall i = 1,\dots,n$;
	\end{itemize}
	Variável de decisão: $x_i \in \{0,1\}$, $\forall i = 1,\dots,n$.
	$$
		\text{maximize}~Z = \sum_{i=1}^{n} b_i x_i
	$$
	sujeito a:
	$$
		\sum_{i=1}^{n} p_i x_i \leq C, \qquad x_i \in \{0,1\},~\forall i=1,\dots,n
	$$
\end{frame}

\begin{frame}[fragile]
	\frametitle{Arquivo \texttt{mochila.mod}}
	\begin{lstlisting}[language=AMPL]
param n > 0;
param C > 0;
param b{1..n};
param p{1..n};

var x{1..n} binary;

maximize Z: sum{i in 1..n} b[i]*x[i];

subject to capacidade: sum{i in 1..n} p[i]*x[i] <= C;
	\end{lstlisting}
\end{frame}

\begin{frame}[fragile]
	\frametitle{Arquivo \texttt{mochila.dat}}
	\begin{lstlisting}[language=AMPL]
param n := 6;
param C := 10;

param : b   p  :=
        1   8  1
        2  10  4
        3  10  3
        4   8  2
        5   7  0.5
        6   4  2;
	\end{lstlisting}
\end{frame}

\begin{frame}[fragile]
	\frametitle{Arquivo \texttt{mochila.run}}
	\begin{lstlisting}[language=AMPL]
model mochila.mod;
data mochila.dat;
option solver highs;
solve;
display Z;
display x;
	\end{lstlisting}
	\vspace{0.3cm}
	No terminal (Windows ou Linux): \texttt{ampl mochila.run}
\end{frame}

\begin{frame}
	\frametitle{Hoje em aula}
	Vamos construir tudo isso \textbf{ao vivo}, a partir de um Colab
	vazio:
	\begin{enumerate}
		\item Criar conta / logar no site do AMPL e ver como obter a
			licença e instalar o \texttt{amplpy};
		\item Escrever o \texttt{.mod}, \texttt{.dat} e \texttt{.run} da
			mochila, célula por célula;
		\item Resolver com HiGHS e interpretar o resultado.
	\end{enumerate}
	\vspace{0.3cm}
	Acompanhem em: \textbf{OPT001\_intro} e \textbf{OPT002\_knapsack\_compact}.
\end{frame}

%------------------------------------------------------------
\section{Problema da Dieta}
%------------------------------------------------------------

\begin{frame}
	\frametitle{Do refeitório à sua tarefa}
	Sua equipe atua no planejamento do refeitório de uma fábrica, junto
	com a nutricionista responsável. É preciso \textbf{reduzir o custo}
	das refeições \textbf{sem prejudicar a saúde} dos trabalhadores.
	\vspace{0.3cm}

	A nutricionista fixou, por dia, quantidades mínimas de:
	\begin{itemize}
		\item Proteínas, cálcio, sódio, ferro e vitaminas;
	\end{itemize}
	e um limite \textbf{máximo} de calorias.
	\vspace{0.3cm}
	Esse é o clássico \textbf{Problema da Dieta}: escolher alimentos e
	quantidades que atendam às exigências nutricionais ao menor custo.
\end{frame}

\begin{frame}
	\frametitle{Problema da Dieta -- semelhança com a mochila}
	\begin{itemize}
		\item Mochila: \textbf{maximiza} benefício sob um limite (peso);
		\item Dieta: \textbf{minimiza} custo sob \textbf{vários} limites
			(um para cada nutriente), com limite inferior e superior;
		\item Mochila: variáveis \textbf{binárias} (leva ou não leva);
		\item Dieta: variáveis \textbf{contínuas} não-negativas (quantidade
			de cada alimento).
	\end{itemize}
	\vspace{0.3cm}
	A estrutura de modelagem (parâmetros, variáveis, objetivo,
	restrições) é a mesma ideia -- muda o que está sendo decidido.
\end{frame}

%------------------------------------------------------------
\section{Tarefa da Aula 1}
%------------------------------------------------------------

\begin{frame}
	\frametitle{Tarefa da Aula 1}
	\begin{enumerate}
		\item Pesquisem, por conta própria, dados nutricionais reais de
			alimentos (tabela TACO, rótulos, bases nutricionais, etc.);
		\item Montem \textbf{seu próprio} \texttt{dieta.dat}, com pelo
			menos 10 alimentos e as propriedades nutricionais pedidas;
		\item Escrevam a formulação compacta do problema (parâmetros,
			variáveis, objetivo, restrições);
		\item Implementem em AMPL (\texttt{.mod}, \texttt{.dat},
			\texttt{.run}) e resolvam com HiGHS;
	\end{enumerate}
	\vspace{0.2cm}
	Detalhes completos em \texttt{tarefas/tarefa01-dieta.md} na página
	da aula.
\end{frame}

\begin{frame}
	\frametitle{Perguntas para reflexão}
	\begin{enumerate}
		\item Quais as limitações da solução obtida?
		\item O que poderia ser melhorado no modelo?
		\item Como pensar o cardápio de um refeitório real (ex.: RU do
			campus), ao longo de uma semana, mês ou ano?
		\item Este modelo poderia ser aplicado na sua casa ou trabalho?
	\end{enumerate}
\end{frame}

\begin{frame}
	\frametitle{Referências}
	\begin{itemize}
		\item Notebooks \texttt{OPT001\_intro} e
			\texttt{OPT002\_knapsack\_compact} (material da disciplina);
		\item AMPL: \url{https://ampl.com};
		\item Dados de referência do problema da dieta: Tirone (2019),
			\url{https://imef.furg.br/images/stories/Monografias/Matematica_aplicada/2019/2019-2_Giulia.pdf}.
	\end{itemize}
\end{frame}

\end{document}
